\subsection{Corte temprano}

Tenemos una trayectoria determinista para el contador y una estimación de
concentración mientras queda una fracción positiva de piezas. Elegimos ahora el
corte de modo que las reinas colocadas después puedan destruir sólo una parte de
las reparaciones ya disponibles. Para ello fijamos
\begin{equation}\label{eq:constantes}
 \eta=\kappa/2,\qquad a_0=\eta^4/16,\qquad \delta_0=a_0/16;
\end{equation}
así $a_0>8\delta_0$ y $\delta_0<1/4$. Para un corte $\delta\in(0,1/4)$ sea
$m_\delta=\lceil\delta n\rceil$, $\tau_\delta$ el instante en que quedan exactamente $m_\delta$
piezas reales pendientes, y $\zeta=\delta/4$. El corte $\delta_0$ se usa en la cota inferior
(Sección~\ref{sec:positiva}); en las Secciones~\ref{sec:inversa} y~\ref{sec:limite} el corte
$\delta$ es libre y se hace tender a cero.

\begin{lema}[Absorbedores en el corte]\label{lem:corte}
Sea $h\in\{1,2\}$ y $\delta=\delta_0$. Salvo un evento de probabilidad $e^{-cn}$ bajo $\QQ_h$, se
alcanza el corte con $\tau_{\delta_0}\le T_{\delta_0/4}$ y, simultáneamente para todas las raíces,
\begin{equation}\label{eq:corte}
 A^{rc}_{\tau_{\delta_0}}\ge a_0n .
\end{equation}
\end{lema}

\begin{proof}
La concentración de los colores en~\eqref{eq:regularidad}, con tolerancia menor que $\delta_0/8$,
asegura $\tau_{\delta_0}\le T_{\delta_0/4}$ salvo un evento de probabilidad $e^{-cn}$. Determinísticamente,
en el corte $y=\tau_{\delta_0}/N\ge(q-m_{\delta_0})/N\ge1/4$ para $n$ grande. Además
$\sum_{e\in\mathcal I_{rc}}\pi_e\ge q-O(1)\ge n/2$, pues $\sum_e\pi_e=q$ y las casillas excluidas
son $O(n)$ con $\pi_e\le C/n$. Por el Lema~\ref{lem:balance}, $1-yp_{d_i}\ge1-p_{d_i}\ge\eta$, de
modo que $a_n^{rc}(y)\ge\eta^4n/8$ para todo $y\ge1/4$. El Lema~\ref{lem:conc_directa} con
tolerancia $\gamma=\eta^4/16$, aplicado a todos los horizontes deterministas $T\le T_{\delta_0/4}$ con
$T/N\ge1/4$ (unión sobre $O(n)$ horizontes, lo que permite el corte aleatorio), da
$A^{rc}_{\tau_{\delta_0}}\ge\eta^4n/16=a_0n$ para todas las raíces. La afirmación sólo requiere que
se alcance el corte; un fallo previo ya está cubierto por el Lema~\ref{lem:regularidad}.
\end{proof}

\emph{Orden de las elecciones.} Se fijan $\eta,a_0,\delta_0$ y $\zeta=\delta_0/4$; después las tolerancias de los
Lemas~\ref{lem:regularidad}--\ref{lem:conc_directa}, que dependen de ellas; finalmente se toma
$n$ grande. No se exige uniformidad cuando $\delta\to0$. Las constantes son analíticas y muy
conservadoras.

\subsection{Pérdida determinista después del corte}

\begin{lema}\label{lem:perdida}
Sea $\delta\in(0,1/4)$ y $(r,c)$ una raíz cuya fila y cuya columna permanecen físicamente vacías. Cada pieza real
crea a lo sumo una candidata ocupada de $(r,c)$ y anula el término de a lo sumo cuatro; las piezas
virtuales no cambian $A^{rc}$. En consecuencia, en el evento de éxito y para todo cruce
finalmente vacío,
\[
 A^{rc}_{\mathrm{final}}\ge A^{rc}_{\tau_\delta}-4m_\delta,\qquad
 |A^{rc}_{\mathrm{final}}-A^{rc}_{\tau_\delta}|\le5m_\delta .
\]
\end{lema}

\begin{proof}
Una nueva reina en la casilla $g$ sólo puede crear el término $Z^\oc_{g,\varnothing}$. Destruye
el término de una candidata ocupada $e$ si $g$ está en una de las diagonales de destino de $e$;
$g$ está en una diagonal y una antidiagonal, y cada una de ellas es diagonal de destino de a lo
sumo dos candidatas ocupadas, por unicidad de filas y columnas. Las piezas virtuales sólo
desactivan vértices, y $A^{rc}$ depende únicamente de $O$ y $V$. Tras el corte se colocan
exactamente $m_\delta$ piezas reales.
\end{proof}
